% =============================================================================
% CHAPTER 2: Background
% =============================================================================
% SPDX-FileCopyrightText: 2026 Frank Langbein <frank@langbein.org>
% SPDX-License-Identifier: LPPL-1.3c
%
% Suggested sections: Context (or Wider context / Domain and context); Related
% work (or Literature review / Prior work); Problem statement (and research
% question(s)). Optional: Theory and concepts; Stakeholders and constraints
% (as subsections of Context or separate). End with a clear problem statement.
% Assume a competent reader; avoid lengthy descriptions of standard tools.
% =============================================================================

\section{Introduction}\label{sec:background-intro}

This chapter reviews the educational and technical background relating to VulnForge. It considers practical cyber security training, the role of \glspl{ctf} and cyber ranges, then examines prior attempts at procedural challenge generation. It then discusses evidence that \glspl{llm} can generate insecure code, the role of adaptive tutoring, and the containment issues related to deploying deliberately vulnerable machines. The chapter concludes by identifying the gap that VulnForge addresses: whether generated, validated, and short-lived vulnerable environments can support reusable cyber security training in a controlled setting.

\section{Context}\label{sec:context}

Cyber security education has steadily moved towards practical, scenario-based learning, where learners develop skills by applying techniques rather than only learning definitions. Students need to understand how vulnerabilities are discovered, exploited, and mitigated in realistic settings, rather than only recognising terms such as SQL injection or privilege escalation. \gls{ctf} exercises have therefore become a widely used educational tool, with prior work arguing that they complement traditional teaching and cover knowledge and skills in ways that map well to cyber security training~\citep{svabensky_cybersecurity_2021}.

Within this context, VulnForge focuses on a specific design problem: whether generative AI can be used to create and manage personalised vulnerable training environments that are useful for practice while remaining operationally safe and practically deployable.

Many existing cyber security training platforms rely on static, pre-authored challenges that expose the learner to the same vulnerability, scenario, and solution path each time. At the time of writing, the platforms reviewed for this project appear to rely primarily on curated or pre-authored content rather than on-demand synthesis of new vulnerable services paired with a grounded tutor.

Practical platforms can help turn abstract concepts into simulated practice. However, many commercial and community platforms are built around a shared architectural assumption: once a challenge is authored, reviewed, and published, users encounter an identical or near-identical artefact. This supports reliability and content quality, but it can also reduce replayability and limit the learner's opportunity to explore variations of the same underlying concept. For example, in a web vulnerability task, a platform may guide the user towards exploiting a path traversal~\citep{noauthor_cwe_nodate}, but the learner may have fewer opportunities to independently identify that vulnerability type across varied scenarios.

The demand for these cyber security skills has grown considerably faster than education around them. The 2024 ISC2 Cybersecurity Workforce study estimated that the global workforce gap in relation to professionals needed to defend organisations reached approximately 4.76 million in 2024, a 19.1 percent increase from 2023~\citep{noauthor_results_nodate}. The study also reported that 30 percent of respondents reported skills gaps within their teams, and 33 percent reported staffing shortages. In addition, the United States Bureau of Labor Statistics projects a 29 percent growth rate for information security analysts through to 2034~\citep{noauthor_information_nodate}. These figures suggest a continuing need for practical training approaches that can scale access to realistic skill development opportunities.

\section{Related work}\label{sec:related}

Existing platforms have demonstrated the value of hands-on training. Hack The Box Academy provides a structured CTF-style route, presenting cyber security training through target hosts, checkpoints, exercises, and skill assessments that allow learners to reproduce techniques they have learned about rather than simply read about them. Katsantonis et al. (2023)~\citep{katsantonis_cyber_2023} argues that cyber ranges have become important learning and training tools, but note that their design is often hindered by a lack of shared standards and methodologies. VulnForge is not only concerned with generating vulnerable code; it also requires a broader system in which an LLM produces structured vulnerability specifications, an orchestration layer deploys these into isolated challenge environments, learners access them remotely, and a tutor provides contextual hints. The proposed solution is therefore closer to a lightweight cyber range than a conventional CTF platform. Lazarov et al. (2025)~\citep{lazarov_lessons_2025} reinforce this distinction by describing cyber ranges as controlled environments for deploying virtual machines and networks, and by arguing that interactivity can improve engagement and practical skill development. For VulnForge, the relevant design requirement is therefore not only to create exploitable targets, but to embed them within a system that can support structured practice.

TryHackMe~\citep{noauthor_tryhackme_nodate}, founded in 2018, takes a beginner-oriented approach, organising content into rooms that combine guided labs, written instructions, and deployable targets. Similar to Hack The Box, many such platforms rely on repeatable challenge artefacts and solution paths, with some modes varying flags or scoring conditions while leaving the underlying vulnerability unchanged. picoCTF~\citep{picoctf} is the largest free student-oriented CTF and also follows a curated challenge model, often cloning vulnerable web applications for learners rather than generating new vulnerability instances on demand.

Including aspects of a CTF is particularly relevant here because it provides a manageable pedagogical unit: a bounded task with a clear goal, observable progress, and explicit proof of successful exploitation. Švábenský et al. (2021)~\citep{svabensky_automated_2024} describe CTFs as a "popular form of cybersecurity education, where students solve hands-on tasks in an informal, game-like setting". However, traditional CTF content raises questions about replayability, variety, and personalisation. Once users have exhausted a fixed content library, the value of repeated practice may decline. Švábenský's discussion of personalised training paths supports the value of adapting challenge selection and sequencing to learners, which is one motivation for VulnForge's generated challenge model.

These platforms and studies show that practical, challenge-based security training is valuable, but also highlight the importance of variation, replayability, and contextual support. Existing platforms provide structured environments in which learners can practise techniques, but much of their value depends on the availability of suitably challenging content. VulnForge builds on this model, but explores whether challenge content can be generated dynamically while still being deployed, validated, and supported through a controlled learning environment.


\section{Procedural Challenge Generation}\label{sec:procedural}
Several systems have attempted to introduce more variety into CTF content. The foundational contribution here is Automatic Problem Generation (APG) by Burket et al. (2015) at the USENIX Summit on Gaming, Games, and Gamification in Security Education~\citep{burket_automatic_2015}. APG was motivated by flag sharing behaviour observed during picoCTF in 2014. This system generates many instances of a challenge from a single parameterised template, ensuring that each player receives a unique flag and, in some tasks, a unique problem structure. However, this approach was not extensible to challenge types that could not be captured by a declarative template. AutoCTF took a different approach, extending the LAVA binary analysis framework to inject memory corruption bugs into existing codebases. A four-university trial demonstrated that students could solve these automatically, but the technique was tightly coupled with LAVA's data flow and did not extend beyond memory corruption bugs in C programs.

The closest precursor to the wide adoption of LLMs was the Security Scenario Generator~\citep{schreuders_security_nodate} (SecGen), a Ruby-based framework that composes randomly selected Puppet modules, each encoding a specific service misconfiguration, known CVE, or user account vulnerability into a virtual machine. This supported a rich variety of scenarios and included a hints system. However, each vulnerability still had to be hand-authored, with the system randomising the selection of Puppet modules rather than the vulnerability itself. APACA~\citep{eckroth_alpaca_2019} extended the idea of a composable environment using AI planning over structured vulnerabilities and configurations to construct multi-chain exploits. However, similar to SecGen, it relied on a pre-populated library.

The most recent example of an LLM-driven agent is AgentCyTE~\citep{rodriguez_agentcyte_2025}. This uses XML files for the CORE network emulator, integrating CVEs from VulnHub into training scenarios. It is the closest current example to the LLM-generated approach that VulnForge explores. However, AgentCyTE uses the LLM to modify scenario creation around a library of pre-selected CVE-backed Docker images; it does not synthesise new vulnerabilities.

\begin{longtable}{|p{0.13\textwidth}|p{0.31\textwidth}|p{0.36\textwidth}|}
	\caption{Comparison of procedural and automated challenge generation systems}
	\label{tab:procedural-generation-comparison} \\
	\hline
	\textbf{System} & \textbf{Generation method}         & \textbf{Limitation relevant to VulnForge}   \\
	\hline
	\endfirsthead
	\hline
	\textbf{System} & \textbf{Generation method}         & \textbf{Limitation relevant to VulnForge}   \\
	\hline
	\endhead
	APG             & Parameterised templates            & Limited to predefined structures            \\
	\hline
	AutoCTF         & Injects bugs into C programs       & Focused on memory corruption                \\
	\hline
	SecGen          & Composes Puppet modules            & Requires hand-authored modules              \\
	\hline
	APACA           & AI planning over known components  & Depends on predefined vulnerability library \\
	\hline
	AgentCyTE       & LLM-assisted scenario construction & Uses pre-selected CVE-backed images         \\
	\hline
\end{longtable}

These systems demonstrate that automated or semi-automated challenge construction is feasible, but they generally depend on predefined templates, hand-authored modules, existing vulnerable images, or narrow vulnerability classes. This leaves a gap for systems that can create new challenge specifications on demand while still remaining constrained enough to deploy safely. VulnForge investigates this design point by using an LLM to generate structured vulnerability specifications dynamically, while relying on schema validation, policy checks, deployment checks, and containment controls to manage the risks of model-generated content.

\section{LLMs as authors of insecure code}\label{sec:LLMs}
Several studies suggest that LLMs can produce security-relevant code, including code containing exploitable weaknesses. A key example is Pearce et al. (2022)~\citep{pearce_asleep_2021}, whose systematic evaluation of GitHub Copilot's code contributions across 89 scenarios drawn from MITRE's 2021 Top 25 Common Weakness Enumeration found that approximately 40 percent of generated programs contained exploitable vulnerabilities. These were concentrated in CWEs such as buffer overflows, SQL injection, and use-after-free. This evidence does not mean that LLMs are reliable challenge authors by default, but it does show that they can produce vulnerability classes that are relevant to security training.

A more concerning dimension of this behaviour was established by Perry et al. (2023)~\citep{perry_users_2023}, whose experiment with 47 participants found that AI-assisted programming produced significantly less secure code while participants simultaneously believed that their code was secure. This suggests that LLMs can not only produce insecure code, but also distort developer confidence in that code. This matters because an LLM-based training platform must not simply produce code and assume correctness. It should treat the generated output as untrusted until validated. In VulnForge, this motivates the use of a structured vulnerability schema, server-side flag generation, deployment checks, and evaluation of whether the generated challenge is actually exploitable.

Offensive agent literature also suggests that LLMs can reason about vulnerabilities and exploitation tasks, although performance varies by task type and setup. PentestGPT~\citep{deng_pentestgpt_2024} achieved a 228.6 percent improvement over a GPT-3.5 baseline on penetration testing benchmarks. Cybench evaluates language model agents on professional CTF tasks and shows that frontier models can solve some easier tasks, especially when provided with scaffolding or subtasks~\citep{zhang_cybench_2025}. NYU CTF Bench similarly evaluates models against CTF-like challenges and reports both capability and limitations, including weaker performance on some challenge types such as incident response~\citep{shao_nyu_2025}. These systems are relevant because VulnForge gives the model two roles: author and tutor. The model creates the vulnerability, but it must also retain enough context to support hints. This influenced the design decision that the LLM should not output arbitrary prose, but instead a schema-constrained vulnerability containing details about services, files, packages, validation steps, and a solution summary.

However, accidental vulnerability generation is not the same as reliable educational challenge generation. An LLM may produce insecure code, but that does not necessarily mean that the vulnerability is intentional, reachable, understandable, or safe to expose to a learner. VulnForge therefore treats LLM output as untrusted: the model is asked to meet an internal specification, and the system performs validation and deployment checks before exposing the result to the user.

\section{LLM tutoring}\label{sec:llm-tutoring}
A tutor is included in this project because practical cyber security tasks can quickly become frustrating when a learner is stuck without feedback. A static hint list can reveal too much, while a full walkthrough undermines the learning objective. An adaptive tutor can provide progressive scaffolding by considering what the user has tried, suggesting enumeration strategies, and explaining relevant concepts when needed.

The idea of a tutor to aid learning is supported in educational psychology. Bloom et al. (1984)~\citep{bloom_2_1984} argued that one-to-one tutoring can produce more significant learning gains than conventional classroom instruction. VanLehn et al. (2011)~\citep{vanlehn_relative_2011} reviewed intelligent tutoring systems and found that well-designed tutoring systems can approach the effectiveness of human tutoring in some cases. Although these results may not directly transfer to CTFs, they support the general educational principle that timely, adaptive feedback can be useful to learners attempting difficult tasks.

Lewis et al. (2020)~\citep{lewis_retrieval-augmented_2021} describe Retrieval Augmented Generation (RAG) as a way to ground responses in context rather than relying only on model memory. VulnForge uses a similar principle: the tutor is grounded in the generated challenge metadata, including the attack vector, hint, and solution summary. This allows it to provide challenge-specific guidance rather than only generic advice about vulnerabilities.

A risk with the tutor is revealing too much. Liffiton et al. (2023)~\citep{liffiton_codehelp_2024} wrote about CodeHelp, a guardrailed LLM for programming education, that gives guidance without providing the solution. VulnForge aims to take a similar approach, it should scaffold reasoning without giving the flag or exact exploit sequence.

Prompt injection is another concern when using a chatbot in an adversarial learning environment. OWASP lists prompt injection in the OWASP Top 10 for LLM Applications~\citep{editor_owasp_nodate}. Wei et al. (2023)~\citep{wei_jailbroken_2023} show that models can fail safety training when instructions conflict or when generalisations do not match the intended policy. In this type of platform, users are adversaries by design: they are expected to bypass restrictions and test assumptions. The implementation therefore needs to combine prompt-level rules, output filtering, application-side checks, and other guardrails. The tutor should be instructed not to share the complete solution, and the application should check for forbidden flag patterns. This does not eliminate prompt injection risk, but it makes the risk an explicit design consideration.

For VulnForge, this means that the tutor cannot be treated as a general-purpose chatbot. It must be grounded in challenge metadata, constrained by rules that prevent solution disclosure, and checked by application-level safeguards that block solution leakage. The tutoring component is therefore part of the system design rather than an optional feature in the interface.

\section{Sandboxed deployment}\label{sec:sandboxed}
Running environments that are deliberately insecure present a host of security problems. The environment has to be vulnerable enough where the learner is able to exploit it, but not so exposed that a successful vulnerability is able to escape the intended boundary and attack other systems on the network or host. This makes the development of VulnForge different to simply a web application, it has to make many considerations for creating targets that enable this adversarial action.

NIST SP 800-190 identifies container images, registries, orchestrators, runtimes, and hosts as distinct areas of container security concern~\citep{souppaya_application_2017}. This presents an important lesson for VulnForge: containers should not be automatically treated as strong security boundaries. Sultan et al. (2019)~\citep{noauthor_pdf_nodate} similarly described container security as a layered problem, involving kernel configuration, namespaces, mandatory access control, and runtime configuration.

This matters because VulnForge uses Proxmox VE to host unprivileged containers, with Ansible handling container lifecycle and vulnerability injection. LXC containers are appropriate for the scale of the project because they are lightweight, scriptable, and integrate well with Proxmox VE. However, they do not provide the same boundaries as regular virtual machines (VMs). Standard VMs do not share a kernel or namespace with the host, providing additional isolation against attacks. A production-ready version of this application would require stronger isolation, more network policies, aggressive egress control, and possibly technologies such as Firecracker~\citep{noauthor_firecracker_nodate}, which is purpose-built for managing secure multi-tenant containers.

VulnForge does not attempt to solve multi-tenant sandboxing at a commercial scale. Instead, it evaluates whether a prototype can combine unprivileged containers, non-root learner accounts, time-limited deployment, controlled access paths, and cleanup mechanisms to reduce risk sufficiently for a controlled educational environment. This makes containment and ephemerality central design requirements rather than secondary deployment details.

Ephemerality is therefore a core part of the deployment model. If a container remains online continually, it becomes a standing target, consumes resources, and increases the host's exposure to attack. Time-limited environments reduce this risk and also act as a resource control mechanism, preventing users from spinning up too many containers and causing excessive load on the service. Network isolation is also important: the containers should not be usable as a pivot point to the wider Internet or as a route to other learners' containers.

\section{Problem statement}\label{sec:problem}

The background discussion identifies a gap in existing practical cyber security training systems. CTFs and cyber ranges provide valuable hands-on learning, but many rely on static or pre-authored challenge artefacts. Prior procedural generation systems improve variety, but often depend on templates, curated modules, or existing vulnerable images. LLMs offer a possible route to more flexible challenge generation, but their outputs are unreliable and potentially unsafe unless they are constrained, validated, and deployed within controlled infrastructure.

VulnForge therefore investigates whether generative AI can move from being an external assistant to becoming part of the content production and delivery pipeline itself. The system aims to generate vulnerability specifications, deploy them into isolated short-lived environments, support the learner with contextual guidance, and clean up resources afterwards.

This leads to the following research questions:
\begin{enumerate}
	\item To what extent can an \gls{llm}-based pipeline generate vulnerability specifications that are valid, deployable, and sufficiently varied within defined challenge categories?
	\item How can generated challenges be deployed and exposed through isolated, short-lived environments without undermining host safety, challenge integrity, or resource control?
	\item What evidence suggests that an on-demand generated approach could improve replayability, variation, and opportunities for guided practice compared with a purely static challenge model?
\end{enumerate}
